\documentclass[11pt]{article}
\usepackage[T1]{fontenc}
\usepackage{lmodern}
\usepackage{amsmath,amssymb,amsthm}
\usepackage{graphicx}
\usepackage{hyperref}
\usepackage[margin=1in]{geometry}

\newtheorem{theorem}{Theorem}
\newtheorem{remark}{Remark}

\newcommand{\G}{\mathcal G}
\newcommand{\HL}{\mathcal HL}
\newcommand{\N}{\mathbb N}
\newcommand{\C}{\mathbb C}
\newcommand{\Id}{\operatorname{Id}}

\title{Strict dilation approximants for the BAP question in holomorphic-Lipschitz linearization}
\author{arXiv automatic math run \texttt{fa\_banach\_001}, lane 10}
\date{June 18, 2026}

\begin{document}
\maketitle

\section*{Status}

This packet records a likely valid partial result related to Question 1 of
Aron--Dimant--Garcia-Lirola--Maestre \cite{ADGM}.  The source asks whether
\(\G_0(B_X)\) has the bounded approximation property whenever \(X\) has the
bounded approximation property.  The result below does not settle the endpoint
identity operator.  It proves that the natural dilation operators
\[
  D_r\delta(x)=\delta(rx)
\]
are bounded-strong limits of finite-rank operators whenever \(r\) lies strictly
below the reciprocal of a BAP constant for \(X\).  Thus the source's scaling
obstruction is sharp at the level of strict dilations.

\section*{Source question}

The question appears after the source proof that \(X\) has the MAP if and only
if \(\G_0(B_X)\) has the MAP.  The authors note that their MAP argument does
not adapt directly to BAP because finite-rank approximations of the identity
may send \(B_X\) into a larger ball.

\begin{figure}[h]
\centering
\includegraphics[width=\textwidth]{figures/open_problem_crop.png}
\caption{Question 1 on page 12 of arXiv:2304.07149, together with the scaling
obstruction immediately preceding it.}
\label{fig:question}
\end{figure}

\section*{Definitions and notation}

Let \(X\) be a complex Banach space and \(B_X\) its open unit ball.  The space
\(\HL_0(B_X)\) consists of scalar holomorphic Lipschitz functions on \(B_X\)
vanishing at \(0\), with norm equal to the Lipschitz constant.  Its canonical
predual is denoted \(\G_0(B_X)\), and \(\delta:B_X\to \G_0(B_X)\) is the
evaluation map.  The source paper proves the linearization property: every
\(\phi\in\HL_0(B_X,Y)\) has a unique bounded operator
\[
  T_\phi:\G_0(B_X)\to Y,\qquad T_\phi\delta(x)=\phi(x),
\]
with \(\|T_\phi\|=L(\phi)\).

If \(X\) has \(\lambda\)-BAP, we fix a net \((A_\alpha)\) of finite-rank
operators on \(X\) such that \(\|A_\alpha\|\leq \lambda\) and
\(A_\alpha x\to x\) for each \(x\in X\).

\section*{Proof intuition}

The source MAP proof fails for BAP because \(A_\alpha(B_X)\) may leave
\(B_X\).  If we first shrink by a factor \(r<1/\lambda\), however,
\(rA_\alpha(B_X)\subset B_X\), so the holomorphic-Lipschitz linearization is
available.  The resulting operator need not be finite rank, since even a
finite-dimensional ball has an infinite-dimensional holomorphic-Lipschitz
free space.  But it factors through \(\G_0(B_E)\), where \(E=A_\alpha X\) is
finite-dimensional; \(\G_0(B_E)\) has the MAP by the source theorem.  Replacing
the middle identity by finite-rank MAP approximants makes the factorization
finite rank without changing the uniform bound.  Letting \(A_\alpha\to I\)
then gives the strict dilation \(D_r\).

\section*{The partial theorem}

\begin{theorem}
Let \(X\) be a complex Banach space with \(\lambda\)-BAP.  For every
\(0<r<1/\lambda\), define \(D_r:\G_0(B_X)\to\G_0(B_X)\) by
\[
  D_r(\delta(x))=\delta(rx),\qquad x\in B_X.
\]
Then \(D_r\) belongs to the strong operator closure of the finite-rank
operators on \(\G_0(B_X)\) of norm at most \(r\lambda\).  In particular, every
strict dilation below the BAP radius has a uniformly bounded finite-rank
approximation scheme on \(\G_0(B_X)\).
\end{theorem}

\begin{proof}
Fix \(r\) with \(0<r<1/\lambda\).  Let
\((A_\alpha)\subset\mathcal L(X,X)\) be finite-rank, \(\|A_\alpha\|\leq\lambda\),
and \(A_\alpha x\to x\) pointwise.  Put \(E_\alpha=A_\alpha X\), equipped with
the inherited norm from \(X\).

The map
\[
  \phi_\alpha:B_X\to B_{E_\alpha},\qquad
  \phi_\alpha(x)=rA_\alpha x
\]
is holomorphic, fixes \(0\), and is Lipschitz with
\[
  L(\phi_\alpha)\leq r\lambda<1.
\]
By linearization it induces
\[
  U_\alpha:\G_0(B_X)\to \G_0(B_{E_\alpha}),\qquad
  U_\alpha\delta_X(x)=\delta_{E_\alpha}(rA_\alpha x),
\]
with \(\|U_\alpha\|\leq r\lambda\).

Let \(j_\alpha:E_\alpha\hookrightarrow X\) be the inclusion.  Since
\(\|j_\alpha\|=1\), the canonical composition/linearization map
\[
  J_\alpha:\G_0(B_{E_\alpha})\to \G_0(B_X),\qquad
  J_\alpha\delta_{E_\alpha}(y)=\delta_X(y),
\]
has norm at most \(1\).  Hence
\[
  S_\alpha:=J_\alpha U_\alpha
\]
satisfies \(\|S_\alpha\|\leq r\lambda\) and
\[
  S_\alpha\delta_X(x)=\delta_X(rA_\alpha x).
\]

For each \(\alpha\), the space \(E_\alpha\) is finite-dimensional, hence has
the MAP.  By Theorem 3.2 of \cite{ADGM}, \(\G_0(B_{E_\alpha})\) has the MAP.
Therefore there is a net \((R_{\alpha,\beta})_\beta\) of finite-rank operators
on \(\G_0(B_{E_\alpha})\), with \(\|R_{\alpha,\beta}\|\leq1\), converging
strongly to the identity on \(\G_0(B_{E_\alpha})\).  The operators
\[
  F_{\alpha,\beta}:=J_\alpha R_{\alpha,\beta}U_\alpha
  :\G_0(B_X)\to\G_0(B_X)
\]
are finite rank and satisfy
\[
  \|F_{\alpha,\beta}\|\leq r\lambda.
\]
For fixed \(\alpha\), they converge strongly to \(S_\alpha\).

Finally, \(S_\alpha\) converges strongly to \(D_r\).  Indeed, for each
\(x\in B_X\),
\[
  S_\alpha\delta_X(x)=\delta_X(rA_\alpha x)\longrightarrow \delta_X(rx)
  =D_r\delta_X(x),
\]
because \(\delta_X\) is Lipschitz and \(A_\alpha x\to x\).  The operators
\((S_\alpha)\) are uniformly bounded by \(r\lambda\), and
\(\operatorname{span}\{\delta_X(x):x\in B_X\}\) is dense in \(\G_0(B_X)\), so
the convergence holds on all of \(\G_0(B_X)\).

Combining the two strong-closure steps, \(D_r\) lies in the strong operator
closure of finite-rank operators of norm at most \(r\lambda\).
\end{proof}

\begin{remark}
If \(X\) has the MAP, one may take \(\lambda=1\) and let \(r\uparrow1\), which
recovers the endpoint approximation of the identity.  For a genuine
\(\lambda>1\) BAP approximation scheme, the theorem stops at
\(r<1/\lambda\); the missing step is precisely a uniform way to cross the
boundary from strict dilations to \(D_1=\Id\).
\end{remark}

\section*{Verification and novelty checks}

No numerical computation is involved.  The proof uses only the source
linearization theorem, the source MAP equivalence for finite-dimensional
spaces, and the definition of \(\lambda\)-BAP.

Local duplicate checks found an earlier 2304.07149 scaling-obstruction
attempt, which recorded the obstruction but did not promote a theorem.  Web
searches on June 18, 2026 for the exact title, for \(\G_0(B_X)\) with BAP, and
for holomorphic Lipschitz bounded approximation found the source arXiv paper
and the related 2026 composition-operator paper, but no full answer to
Question 1.  The 2026 paper uses BAP for composition-operator
characterizations; it does not appear to settle the BAP of \(\G_0(B_X)\).

\section*{Human-review recommendation}

Review as a modest but clean partial theorem.  The main point to check is the
factorization
\[
  \G_0(B_X)\xrightarrow{U_\alpha}\G_0(B_{E_\alpha})
  \xrightarrow{J_\alpha}\G_0(B_X)
\]
and the use of the MAP of \(\G_0(B_{E_\alpha})\) to replace the middle identity
by finite-rank maps while preserving the \(r\lambda\) bound.

\begin{thebibliography}{9}

\bibitem{ADGM}
R. Aron, V. Dimant, L. C. Garcia-Lirola, and M. Maestre,
\emph{Linearization of holomorphic Lipschitz functions},
arXiv:2304.07149, 2023.

\end{thebibliography}

\end{document}
